\documentclass[aspectratio=169]{beamer}
\input{header}
\coursecode{JKSSB}

\title{Storage Devices}
\subtitle{Lesson 10 --- Magnetic, Optical, Flash and Cloud}
\author{JKSSB Computer Awareness}
\date{\today}

\begin{document}
\frame{\titlepage}

\begin{frame}{Storage in the Memory Hierarchy}
  \begin{itemize}
    \item \key{Primary storage}: directly accessible by the CPU; holds data
      and instructions currently in use (RAM, ROM, cache).
    \item \key{Secondary storage}: non-volatile; keeps data when the power is
      off (hard disk, SSD, optical disc, USB drive).
    \item \key{Tertiary storage}: archival level, often robotic tape
      libraries that are loaded on demand.
  \end{itemize}
  \begin{block}{Anchor point}
    The CPU works with main memory; secondary storage feeds main memory
    through the operating system.
  \end{block}
\end{frame}

\begin{frame}{Magnetic Storage: Hard Disk Drive}
  \begin{itemize}
    \item Data is stored magnetically on \key{rotating platters}.
    \item A \key{read/write head} on an actuator arm moves over the platter
      surface to read or write.
    \item It is a mechanical device: the platters spin and the head moves.
    \item Large capacity at low cost per gigabyte.
  \end{itemize}
  \begin{alertblock}{Remember}
    An HDD is magnetic and mechanical --- it has moving parts.
  \end{alertblock}
\end{frame}

\begin{frame}{Magnetic Tape: Sequential Storage}
  \begin{itemize}
    \item Data is stored magnetically on a long tape in a reel or cartridge.
    \item The tape is read or written in order, one position after another.
    \item This is \key{sequential access}, so reaching a distant record takes
      longer.
    \item Widely used for backup and long-term archival rather than
      day-to-day access.
  \end{itemize}
  \vspace{0.3cm}
  \muted{Tape is cheap and dense, but slow for random lookups.}
\end{frame}

\begin{frame}{Optical Storage}
  \begin{itemize}
    \item Data is stored as microscopic \key{pits and lands} and read by a
      \key{laser}.
    \item Common forms: \key{CD}, \key{DVD}, and \key{Blu-ray} disc.
    \item Optical media are portable, removable, and non-volatile.
    \item Blu-ray uses a shorter-wavelength laser and tighter tracks, so it
      holds more than a DVD.
  \end{itemize}
  \begin{block}{Capacity cue}
    CD $\approx$ 700 MB; DVD $\approx$ 4.7 GB (single layer); Blu-ray
    $\approx$ 25 GB (single layer).
  \end{block}
\end{frame}

\begin{frame}{Optical Media Capacities}
  \begin{center}
    \begin{tabular}{ll}
      \toprule
      \textbf{Medium} & \textbf{Approximate capacity (single layer)} \\
      \midrule
      CD & 700 MB \\
      DVD & 4.7 GB \\
      Blu-ray & 25 GB \\
      \bottomrule
    \end{tabular}
  \end{center}
  \vspace{0.3cm}
  \muted{Order from smallest to largest: CD $<$ DVD $<$ Blu-ray.}
\end{frame}

\begin{frame}{Flash and Solid-State Storage}
  \begin{itemize}
    \item \key{Flash memory} stores data in semiconductor cells with
      \key{no moving parts}.
    \item Examples: \key{SSD}, \key{USB flash drive}, and \key{memory card}.
    \item Because nothing spins, flash is faster and more shock-resistant
      than a mechanical drive.
    \item It is non-volatile: data remains when the power is off.
  \end{itemize}
  \begin{alertblock}{High-yield fact}
    Among these devices, the \key{SSD} has the highest transfer rate
    (KB PYQ-27).
  \end{alertblock}
\end{frame}

\begin{frame}{Cloud Storage}
  \begin{itemize}
    \item Data is stored on \key{remote servers} and accessed over the
      Internet.
    \item The provider manages the hardware; the user reaches the files
      through a network connection.
    \item It still ultimately rests on physical secondary storage in a data
      centre.
    \item Useful for backup, sharing, and access from many devices.
  \end{itemize}
  \muted{Cloud storage is remote secondary storage, not main memory.}
\end{frame}

\begin{frame}{HDD versus SSD}
  \begin{columns}[T]
    \column{0.46\textwidth}
      \textbf{Hard disk drive}
      \begin{itemize}
        \item Magnetic storage
        \item Mechanical, moving parts
        \item Rotating platters and a moving head
        \item Usually slower transfer rate
        \item More vulnerable to shock
      \end{itemize}
    \column{0.46\textwidth}
      \textbf{Solid-state drive}
      \begin{itemize}
        \item Flash storage
        \item No moving parts
        \item Semiconductor cells
        \item Higher transfer rate (fastest here)
        \item More shock-resistant
      \end{itemize}
  \end{columns}
\end{frame}

\begin{frame}{Random Access versus Sequential Access}
  \begin{itemize}
    \item \key{Random (direct) access}: any location can be reached in
      approximately equal time. RAM and hard disks behave this way.
    \item \key{Sequential access}: locations are reached in order; the time
      depends on how far the target is from the current position. Magnetic
      tape is sequential.
    \item This distinction explains why tape is poor for direct lookups but
      fine for ordered backup.
  \end{itemize}
  \begin{block}{Quick check}
    Memory and disks are approximately random access; tape is sequential.
  \end{block}
\end{frame}

\begin{frame}{Non-Volatility of Storage}
  \begin{itemize}
    \item \key{Non-volatile} means the contents survive when the power is
      switched off.
    \item Primary storage such as RAM is \key{volatile}: its contents are
      lost without power.
    \item Secondary storage --- HDD, SSD, optical, USB --- is non-volatile,
      so files persist after shutdown.
    \item Do not say RAM is non-volatile; do not say storage is volatile.
  \end{itemize}
  \begin{alertblock}{One-line rule}
    RAM is volatile; storage is non-volatile.
  \end{alertblock}
\end{frame}

\begin{frame}{Office Desktop Example (EX-01)}
  \begin{itemize}
    \item An office desktop uses an \key{SSD} as the system drive, so the
      operating system boots quickly.
    \item A large \key{HDD} holds bulk documents and media at lower cost.
    \item A \key{DVD drive} and \key{USB flash drive} move files between
      machines.
    \item Important files are also copied to \key{cloud storage} for remote
      backup.
  \end{itemize}
  \muted{Fast flash for daily work; magnetic capacity for bulk; optical and
  USB for transfer; cloud for off-site backup.}
\end{frame}

\begin{frame}{Exam Traps to Avoid}
  \begin{itemize}
    \item The \key{SSD} has the highest transfer rate, not the HDD
      (KB PYQ-27).
    \item An HDD is \key{magnetic and mechanical}; an SSD is \key{flash with
      no moving parts}.
    \item Optical capacities run CD $<$ DVD $<$ Blu-ray.
    \item RAM is \key{volatile} but storage is \key{non-volatile}.
  \end{itemize}
  \begin{alertblock}{Exam method}
    Match the storage type to its technology and its access method before
    answering.
  \end{alertblock}
\end{frame}

\begin{frame}{Lesson 10: Recall}
  \begin{itemize}
    \item Primary storage is CPU-accessible; secondary storage is
      non-volatile; tertiary storage is archival.
    \item Magnetic: HDD (rotating platters, moving head) and tape
      (sequential).
    \item Optical: CD $\approx$ 700 MB, DVD $\approx$ 4.7 GB, Blu-ray
      $\approx$ 25 GB (single layer).
    \item Flash: SSD, USB drive, memory card --- no moving parts; the SSD
      has the highest transfer rate.
    \item RAM and disks are approximately random access; tape is sequential.
  \end{itemize}
\end{frame}
\end{document}
